\newpage
\subsubsection{Herramientas para la aplicación web}

La aplicación web se desarrollará como una Single Page Application
(SPA) orientada a la visualización y administración de los datos
sincronizados desde los dispositivos móviles. El stack seleccionado
prioriza la integración nativa entre herramientas, la mantenibilidad
del código y la capacidad de representar datos geoespaciales y
estadísticos de forma clara.

\textbf{Angular} es el framework principal para el desarrollo del
frontend. Desarrollado y mantenido por Google, Angular ofrece una
arquitectura basada en componentes, inyección de dependencias nativa
y un sistema de módulos que facilita la organización del código en
aplicaciones de mediana y gran escala. Su integración con TypeScript
garantiza tipado estático y detección temprana de errores durante
el desarrollo.

\textbf{Angular Material} es la biblioteca oficial de componentes
de interfaz de usuario para Angular, basada en los principios de
Material Design. Proveerá los componentes visuales del sistema como
tablas, formularios, diálogos, navegación y tarjetas, asegurando
una interfaz consistente, accesible y alineada con los estándares
modernos de diseño.

\textbf{D3.js} se utilizará para la visualización de datos estadísticos y geoespaciales del sistema. Su capacidad para construir visualizaciones personalizadas basadas en datos permitirá representar la distribución de focos de infestación, la evolución temporal del estado sanitario del arbolado y las estadísticas generales del sistema de forma interactiva y precisa.

\textbf{MapLibre GL JS} es la biblioteca de código abierto utilizada
para la visualización de mapas vectoriales en la aplicación web.
Permite renderizar mapas interactivos directamente en el navegador
y representar mediante pines georreferenciados las ubicaciones de
los registros sincronizados desde los dispositivos móviles, facilitando la supervisión espacial del estado sanitario del arbolado desde la plataforma de administración.

\textbf{HttpClient} es el módulo nativo de Angular para la gestión
de peticiones HTTP. Se utilizará para consumir los endpoints REST
expuestos por Amazon API Gateway, gestionando la comunicación entre
el frontend y el backend de forma eficiente mediante observables y
soporte para interceptores, los cuales permitirán adjuntar
automáticamente el token de autenticación de Amazon Cognito en cada
petición.

\textbf{Amazon Cognito} gestionará la autenticación de los usuarios
de la plataforma web, incluyendo los roles de Experto de campo y
Administrador. Proveerá el flujo de inicio de sesión, recuperación
de contraseña y gestión de tokens JWT, eliminando la necesidad de
implementar estos mecanismos de forma manual en el backend.

La seguridad en las comunicaciones entre el cliente web y los servicios en la nube se garantizará mediante \textbf{AWS Certificate Manager (ACM)}, el servicio de Amazon que provee y renueva automáticamente certificados SSL/TLS sin costo adicional. ACM se integrará con Amazon CloudFront y Amazon API Gateway para asegurar que toda la comunicación se realice a través de HTTPS, cumpliendo con el requerimiento no funcional RNF-05.